% !TEX root = ../notes.tex

\documentclass[../notes.tex]{subfiles}

\begin{document}

\section{May 6}
Here we go.

\subsection{More on the \texorpdfstring{$p$}{p}-Curvature}
Let's compute with an example.
\begin{remark}
	With $\mc E=\OO$ (or a line bundle more generally), we have $\op{End}\mc E=\OO$, so one can take $\nabla=d+\omega$, where $\omega$ is some $1$-form. Then the $p$-curvature is $\op{Fr}^*(\omega)$, viewed as a section of $\op{Fr}^*\Omega^1$. Explicitly, on vector fields $v$,
	\[\op{curv}_pv=\langle\omega,v\rangle^p.\]
\end{remark}
To compute this, we use the Cartier isomorphism $C$, which is an isomorphism from closed $1$-forms modulo exact forms to $\Omega^1_{X^{(1)}}$.
\begin{example}
	Take $X=\AA^1$. All $1$-forms are closed, and we can compute $\coker d$ has basis $\left\{t^{pi-1}\,dt\right\}_i$. Then the Cartier isomorphism sends $t^{pi-1}\,dt\mapsto T^{i-1}\,dT$; intuitively, one has $C(dt/t)=dT/T$. Now, if $\nabla=d+dt$, one has $\nabla_{\del_t}=\frac d{dt}+1$, and one can compute
	\[\op{curv}_p(d+dt)=dT.\]
	Similarly, one can compute $\op{curv}_p(d+dt/t)=0$. 
\end{example}
\begin{remark}
	In general, using the functoriality from the previous example, one finds $\op{curv}_p(d+df)=df^{(1)}$ and $\op{curv}(d+df/f)=0$, so one can compare this to get
	\[\op{curv}_p(d+\alpha)=\alpha^{(1)}-C(\alpha).\]
	Here, $\alpha^{(1)}$ refers to basically the same object, where we use $T$ instead of $t$. In coordinates, this means $\op{curv}_p(d+a\,dt)=\left(a^p+a^{(p-1)}\right)\,dT$, where $a^{(p-1)}$ means an iterated derivative.
\end{remark}
\begin{example}
	The same sort of calculation shows that $\op{curv}_p(d+dt)=dT$ in general, where $t$ is the local coordinate. Similarly,
	\[\op{curv}_p\left(d+\lambda\frac{dt}t\right)=\left(\lambda^p-\lambda\right)\frac{dT}T.\]
	Notably, this vanishes if and only if $\lambda\pmod p$ lives in $\FF_p$.
\end{example}
\begin{example}
	Suppose $X$ is an elliptic curve over $\QQ$, so we have a unique nowhere vanishing holomorphic $1$-fom $\omega$. It turns out that $\omega_p-C(\omega_p)$ vanishes if and only if $p\mid\#X(\FF_p)$. (By the Weil conjectures, this means that $\#X(\FF_p)=p$ for large $p$.) Accordingly, pairs $(\OO,d+\omega_p)$ with vanishing $p$-curavture are in bijection with line bundles $\mc L$ with $\op{Fr}^*\mc L\cong\OO$.
\end{example}
We now take $\mc E$ to be a vector bundle of higher rank. Here, one cannot view $\op{curv}_p(\mc E,\nabla)$ as an object over $X^{(1)}$. Instead, one needs to choose some conjugation-invariant polynomials of the induced matrix. (Namely, we get a matrix by viewing the curvature as basically outputting to $\op{End}(\mc E)$, with a twist.)
\begin{example}
	Let $(\mc E,\nabla)$ be a vector bundle with flat connection. Then
	\[\tr\op{curv}_p(\mc E,\nabla)^r\]
	is a global section of $\op{Sym}^r\Omega^1_{X^{(1)}}$.
\end{example}
Now is as good a time as any to define the $p$-support, which was already featured last class.
\begin{definition}[support]
	Fix a smooth algebraic variety $X$ over a field $k$ of positive characteristic $p$. Let $(\mc E,\nabla)$ be a vector bundle with flat connection of rank $n$. Then the polynomials $\tr\op{curv}_p(\mc E,\nabla)^r$ cut out a spectral curve $\widetilde X^{(1)}\subseteq T^*X^{(1)}$, which admits a natural map $\widetilde X^{(1)}\to X$ which is finite of degree $n$. Then the \textit{$p$-support} is the subscheme $\op{supp}(\mc E,\nabla)$ of $\widetilde X^{(1)}$.
\end{definition}
\begin{remark}
	Intuitively, over each $x\in X$, we see that $\widetilde X^{(1)}$ simply chooses the $n$ eigenvalues of the corresponding matrix, which explains why the map $\widetilde X^{(1)}\to X$ is finite of degree $n$. This explains why it was called a spectral variety.
\end{remark}
\begin{remark}
	The $p$-support is a module for $Z(D_X)$.
\end{remark}
It is in general hard to compute the $p$-curvature in higher-rank cases, but we can get some partial information sometimes.
\begin{remark}
	Fix $(\mc E,\nabla)$, and suppose that the spectral variety $\widetilde X\subseteq T^*X$ is smooth. Recall that $D_X$ is an Azumaya algebra on $T^*X^{(1)}$, and we let $\theta$ be its class in $\op{Br}T^*X^{(1)}$. It turns out that the following collections are in natural bijection.
	\begin{itemize}
		\item Pairs $(\mc E,\nabla)$ with spectral variety $\widetilde X$.
		\item Line bundles $\mc L$ on $\widetilde X$ for which $\op{curv}_p\nabla_{\mc L}=\theta$.
		\item A splitting of $D_X|_{\widetilde X^{(1)}}$.
	\end{itemize}
	The moral is that passing to the spectral cover allows us to turn vector bundles into line bundles.
\end{remark}
\begin{example}
	Let $X$ be a complete curve of genus bigger than $1$ with canonical line bundle $\mc K$. It turns out that there is a unique non-split extension
	\[0\to\mc K^{1/2}\to\mc E\to\mc K^{-1/2}\to0,\]
	given by the generator in $\mathrm H^1(K)$. Flat connections on $\mc E$ are the same as $\mf{sl}_2$-opers, and they turn out to exist. Furthermore, we will ask for this flat connection to have trivial determinant. Now, for a fixed connection $\nabla_0$, one can write $\nabla=\nabla_0+A$, where $A=\begin{bsmallmatrix}
		0 & \beta \\ 0 & 0
	\end{bsmallmatrix}$, where $\beta$ is a map $\mc K\to\op{End}\mc E$, which has the same data as a global section of $\mc K^2$. On the other hand, $\widetilde X\subseteq T^*X$ is cut out by a section $\alpha$ of $\mc K^2$. It turns out that the map $\beta\mapsto\alpha$ is asymptotic to Frobenius, meaning that the leading term of this polynomial is given by the Frobenius. The idea of the proof is to write $\nabla_t=\nabla_0+t\alpha$, and then we need to compute that $\nabla_t^p$ give $t^p\alpha^p$, which is isospectral to $\alpha^{(1)}$. More details are found in a paper of Chen--Zhu--Bezrukavnikov--Travkin.
\end{example}
% \begin{example}
% 	and $\omega=dt$, one has $C(\omega)=0$. Similarly, taking $X=\mathbb G_m$ and $\omega=\lambda\,dt/t$, then $C(\omega)=(\lambda^p-\lambda)\,dT/T$, where $T$ is the local coordinate on $X^{(1)}$.
% \end{example}

\subsection{Katz's Theorem}
Recall from \Cref{ex:calculus-diagram} that a smooth morphism $f\colon Y\to X$ gives rise to a diagram
% https://q.uiver.app/#q=WzAsMyxbMCwwLCJUXiooWFxcdGltZXNfWFkpIl0sWzEsMCwiVF4qWSJdLFswLDEsIlReKlgiXSxbMCwxLCJkZl4qIiwwLHsic3R5bGUiOnsidGFpbCI6eyJuYW1lIjoiaG9vayIsInNpZGUiOiJ0b3AifX19XSxbMCwyLCJcXG9we3ByfV8xIiwyXV0=&macro_url=https%3A%2F%2Fraw.githubusercontent.com%2FdFoiler%2Fnotes%2Fmaster%2Fnir.tex
\[\begin{tikzcd}[cramped]
	{T^*(X\times_XY)} & {T^*Y} \\
	{T^*X}
	\arrow["{df^*}", hook, from=1-1, to=1-2]
	\arrow["{\op{pr}_1}"', from=1-1, to=2-1]
\end{tikzcd}\]
related to the Gauss--Manin connection. Indeed, let $f_+$ denote the pushforward construction for $\mc D$-modules. Then we saw that the pushforward of $\mc D$-modules by $f_+$ provided an $\hbar$-deformation of $\pi_*df^*$. In positive characteristic, there is another way to find $f_+$ in this diagram.
\begin{lemma} \label{lem:calculus-positive-char}
	Fix a smooth morphism $f\colon Y\to X$ of smooth algebraic varieties of positive characteristic $p$. Then $df^{(1)*}\mc D_Y$ is canonically Morita-equivalent to $\pi^{(1)}\mc D_X$.
\end{lemma}
\begin{proof}
	Omitted.
\end{proof}
\begin{remark}
	Thus, up to the canonical isomorphism in the lemma, we have a composite $\pi_*^{(1)}\circ df^{(1)*}$ which approximately sends $\mc D_Y$-modules to $\mc D_X$-modules. It turns out that this functor is naturally isomorphic to $f_+$.
\end{remark}
If we further assume that $f$ is proper, then we can spread out everything in sight. Unwinding everything, one can show the following.
\begin{proposition}
	Fix a smooth proper algebraic variety $X$. Then the $p$-curvature of the Gauss--Manin connection is nilpotent.
\end{proposition}
\begin{proof}[Idea]
	I think the idea is that the pushforward (via the above description) is supported on the zero section (perhaps seen by passing to characteristic zero in a spreading out argument), but I didn't follow this argument.
\end{proof}
\begin{definition}[Cartier transform]
	Suppose that $X$ lifts to $W_2(k)$ and that $p<n$. Then it turns out that the Azumaya algebra $\mc D_X$ over $\OO_{T^*X^{(1)}}$ splits canonically (as an Azumaya algebra) on a $p$th infinitesimal neighborhood of the zero section. This splitting allows one to use Morita equivalence to see that pairs $(\mc E,\nabla)$ of a vector bundle with flat connection and nilpotent $p$-curvature embed into coherent sheaves on $T^*X^{(1)}$.  In the sequel, we will call this the \textit{Cartier transform}.
\end{definition}
\begin{remark}
	The Cartier transform is compatible with the canonical isomorphism of \Cref{lem:calculus-positive-char}. One can use this to calculate that the Cartier transform of the Gauss--Manin connection is the coherent sheaf
	\[\mathrm H^i\left(\pi_*df^*(\OO_{Y^{(1)}})\right).\]
	Notably, we have already seen this object in the associated graded for the Hodge filtration!
\end{remark}
We are now ready to sketch the following result.
\begin{theorem}[Katz] \label{thm:katz-p-curv}
	\Cref{conj:p-curv} holds for Gauss--Manin connections associated to smooth proper morphisms $Y\to X$.
\end{theorem}
\begin{proof}[Idea]
	Let $\mc E_i$ be the Gauss--Manin connection. If $\op{curv}_p\mc E_i=0$ for infinitely many $p$, then the terms $\op{gr}_{\mathrm{Hodge}}\mc E_i$ are supported on the zero section of $T^*X_\CC$. This means that the Higgs field vanishes, so the terms of the Hodge filtration are $\nabla$-invariant and in particular $\nabla$-flat.

	Now, fix a relatively ample line bundle $\mc L$ on $Y$. This gives a Hermitian, $\nabla$-flat pairing $\langle-,-\rangle$ on $\mc E_i$. In particular, this pairing makes
	\[(-1)^i\langle-,-\rangle|_{\op{gr}\mc E_i}\]
	is positive. Notably, the filtration splits by the previous paragraph, so we can change basis (while remaining $\nabla$-flat) to get a Hermitian positive-definite pairing. Thus, the monodromy lies in a compact (unitary) group, and finiteness follows because monodromy is always a discrete group.
\end{proof}

\end{document}